\section{Results}
\label{sec:results}

\subsection{What agents claim, and how they back it}
\label{sec:claims}
Of the \nClaims runtime-performance PRs, \nNoMech (\pctNoMech) name no
mechanism at all. The most common mechanisms are skipping unneeded work
(\mechSkipUnneededWork), reducing memory use or copies (\mechLessMemory)
and caching results (\mechCacheResults) (Table~\ref{tab:oracle}). Claims
are rarely measured: \evReasoning PRs support the claim by reasoning only
and \evNone offer no evidence (together \pctReasonNone); \evBenchmark cite
a benchmark (\pctBenchmark), \evTest a test, and \evProfile a profile. Only
\nEffect (\pctEffect) state an effect size. \nMeasurableClaims PRs claim
at least one mechanism with a kernel or hardware signal; the rest claim
only GPU, external-API, user-interface or uncategorized mechanisms, or
none.

\subsection{Building the PRs}
\label{sec:funnel}
Of the \nMeasurableClaims PRs, \nGhError are no longer available on GitHub
and \nNoPyFinal do not change Python code in their final diff (AIDev's
per-commit file lists include commits merged in from the base branch),
leaving \nJobs build jobs in one fixed random order. In the first batch
(the first \nAttempted of that order), the build agent produced a
workload that passed the harness checks for \funPassed of \nAttempted
attempted PRs (\pctPassed), at about \usdPerBuild per PR. The others need
an external service (\funNeedsExternalService), add only new code that the
base version cannot run (\funNewCodeOnly), need a GPU (\funNeedsGpu), failed
to install (\funInstallFailed), had unreachable changed code
(\funChangeUnreachable), or did not change Python in a runnable way
(\funNotPythonChange); for \funFailedCheck the agent reported success but
the independent check failed, and \funOther could not be classified.

\subsection{Do the claims hold below the stopwatch?}
\label{sec:measure}
We analyse the \nAnalysed measured PRs (from \nReposAnalysed repositories)
whose claim names at least one checkable mechanism and whose change is
mostly Python; \nLowPy measured PRs whose change is mostly non-Python code
are set aside, because their claim may concern code we do not execute.
The manual review excludes \nExcludedReview PRs: in three the claimed
mechanism lies in code the workload never runs, and in one the base
version fails on the measured path. \nExcludedNoisy PR is excluded because its A/A control (base against base)
itself differs by more than 5\%; for the others the A/A wall-time ratio
stays within \aaMin--\aaMax, so the noise floor is small. A PR is \emph{faster} if its
median wall time drops by at least 5\% and every head run is faster than
every base run; a mechanism \emph{holds} if one of its signals
(Table~\ref{tab:oracle}) moves by at least 10\% in the claimed direction
with complete separation between head and base runs.

\begin{table}[t]
\caption{First measurements: wall time against the claimed mechanism, for
the \nAnalysed analysed PRs.}
\label{tab:quadrant}
\centering\footnotesize
\begin{tabular}{@{}lccc@{}}
\toprule
 & faster & no clear change & slower \\
\midrule
claimed mechanism holds & \qFH & \qNH & \qSH \\
claimed mechanism absent & \qFN & \qNN & \qSN \\
\bottomrule
\end{tabular}
\end{table}

\begin{figure}[t]
\centering
\begin{tikzpicture}
\begin{semilogxaxis}[width=\columnwidth, height=4.2cm,
  xlabel={wall time, head / base (log scale)}, ylabel={PR (sorted)},
  ytick=\empty, xmin=0.03, xmax=15, enlarge y limits=0.07,
  xtick={0.05,0.1,0.2,0.5,1,2,5,10}, xticklabels={0.05,0.1,0.2,0.5,1,2,5,10},
  legend style={font=\scriptsize, at={(0.98,0.04)}, anchor=south east},
  label style={font=\footnotesize}, tick label style={font=\footnotesize}]
\fill[gray!20] ({axis cs:0.95,1}|-{rel axis cs:0,0}) rectangle ({axis cs:1.05,1}|-{rel axis cs:0,1});
\draw[gray] ({axis cs:1,1}|-{rel axis cs:0,0}) -- ({axis cs:1,1}|-{rel axis cs:0,1});
\addplot[only marks, mark=*, black] table[x=ratio, y=rank,
  restrict expr to domain={\thisrow{held}}{1:1}] {fig_wall.dat};
\addplot[only marks, mark=o, black] table[x=ratio, y=rank,
  restrict expr to domain={\thisrow{held}}{0:0}] {fig_wall.dat};
\legend{mechanism holds, mechanism absent}
\end{semilogxaxis}
\end{tikzpicture}
\caption{Wall-time ratio of each analysed PR. The grey band is the
$\pm$5\% zone counted as no clear change.}
\label{fig:wall}
\end{figure}

Table~\ref{tab:quadrant} and Figure~\ref{fig:wall} show the first results. Of the \nAnalysed PRs,
\nFaster are faster, \nNoChange show no clear change and \nSlower are
slower than before the change; the claimed mechanism holds in \nHeld.
When a PR is faster, its claimed mechanism usually appears (\qFH of
\nFaster); when it is slower, it rarely does (\qSH of \nSlower). In this
sample, a missing footprint rarely comes with a speed-up. The rule does not
fire by chance: applied to base against base again, it moves none of
\aaSigTests signal tests (\aaSigRate). It is, however, only partly
specific, because several mechanisms share the instruction count: for the
faster PRs, the signals of mechanisms a PR did \emph{not} claim also move
in \spUnclaimedRateFast of cases, against \spClaimedRateFast for the
claimed ones. A held mechanism is therefore evidence consistent with the
claim, not proof of it.
\nBehaviour PRs also change the output of the workload (for example new
defaults or error messages), as Qi et al.\ also observe for merged agent
fixes~\cite{qi2026merged}. The pattern is the same with and without such
changes: of the \nSame PRs whose output is unchanged, \qFHSame are faster
with the mechanism present and \qFNSame faster without it. \nHiddenCost
PRs make at least one resource clearly worse (system time, memory, page faults, syscalls,
instructions, CPU time or bytes read).

% Case studies. Every number is a macro from analysis/p7_cases.py, which
% reads the medians of the measured runs and asserts full separation.
\paragraph{Faster, but not for the claimed reason}
iPhotron PR~\#453 (Copilot) claims to remove the slowdown of deep SQL
pagination and to load assets on demand instead of all at once. The workload
runs \caseAWall\ of the base wall time (\caseAWallBase\,s to
\caseAWallHead\,s; all ratios use unrounded medians), yet the kernel shows more work, not less: \caseASyscalls\
the syscalls (\caseASyscallsBase\ to \caseASyscallsHead), \caseAPread\ the
\texttt{pread64} calls, \caseABytes\ the bytes read and \caseASys\ the system
time. Total CPU time grows \caseATask\ while the process runs
\caseAThreadsHead\ threads instead of \caseAThreadsBase. Peak memory does not
change. A stopwatch would credit the claim; the trace does not support it.

\paragraph{Slower, and the trace shows where}
iPhotron PR~\#564 (Jules) uses pre-generated thumbnails. On our workload,
which keeps 500 unrelated files in the thumbnail folder, it is
\caseBWall\ slower (\caseBWallBase\,s to \caseBWallHead\,s) and its output
changes. The trace shows a likely cause: \texttt{getdents64}, the call that
lists a directory, rises from \caseBGetdentsBase\ to \caseBGetdentsHead\
(\caseBGetdents), and instructions grow \caseBInstr.

\paragraph{Half the mechanism happens, and it hurts}
A Claude Code PR to gemini-notebook-mcp-cli (\#16) claims non-blocking
I/O and HTTP connection pooling. The first part appears: the main thread
spends \caseDOffcpu\ of its former time off the CPU. The second does not:
the workload opens \caseDSocketHead\ sockets and \caseDConnectHead\
connections on both versions. Meanwhile syscalls grow \caseDSyscalls,
\texttt{futex} calls \caseDFutexc\ and CPU time \caseDTask, and the
workload is \caseDWall\ slower (\caseDWallBase\,s to \caseDWallHead\,s).

\paragraph{A claim that holds}
mlflow PR~\#19520 (Copilot) changes artifact transfer to chunked streaming.
The footprint is exactly the claimed one: \texttt{read} calls fall from
\caseCReadBase\ to \caseCReadHead\ (\caseCRead) for the same bytes read,
instructions fall to \caseCInstr\ and wall time to \caseCWall\ of the base.

